% MANUSCRIPT-IDENTITY: not-active | paper_v | active source is MATH/paper_details/paper_V_high_u_gkba.tex | agent_guidance/paper-registry.yaml
%
% ============================================================================
% Draft replacement text for the two red placeholder slots of the "Numerical
% plan" section of MATH/paper_details/paper_V_high_u_gkba.tex, plus the short
% results section they imply.  Written to be read first and inserted second:
% nothing here has been put into the manuscript, and inserting it is the
% user's decision.
%
% WHAT THE NUMBERS BELOW ARE.  They are exploratory local measurements taken
% on branch claude/paper-v-routes-20260914 between 2026-09-14 and 2026-09-18
% and consolidated in
%   MATH/paper_facing/paper_V_high_u_gkba/cheap_repair_route_results_20260918.md
% (called "the results summary" in the comments below), with the long-horizon
% adaptive-integrator figures used in Block B taken from
%   agent_guidance/paper-v/overnight/joint-gram-t20-20260915.md
% (called "the long-horizon record").  Each number carries its definition, its
% configuration and its horizon where it is used.
%
% WHAT THEY ARE NOT.  They are not adopted evidence.  Every rollout that
% produced them writes classification = exploratory_local_not_promoted; the
% plan of record lists this route as "adoption pending"; the t = 20 entries of
% the standing configuration are one run with no band of their own; and
% adoption, promotion and demotion remain the user's decisions.  For that
% reason every quantitative claim below is wrapped in the draft's own
% \tentative macro, which is what the manuscript already uses to mark a figure
% that has not been adopted.  A band in parentheses is a measured one-setting
% sensitivity, not a statistical error bar.
%
% ONE CONSEQUENTIAL SIDE EFFECT OF INSERTING BLOCK A.  The abstract currently
% ends "No stability, complexity, or quantum-advantage result is claimed in
% this draft."  Block A reports stability diagnostics.  That sentence should
% be re-read before Block A is inserted; it is not edited here.
%
% Three blocks, in the order they would be inserted:
%   Block A  a new \section, to go immediately before \section{Numerical plan}
%   Block B  replacement for the \placeholder under
%            \paragraph*{Regularization slot.}
%   Block C  replacement for the \placeholder under
%            \paragraph*{Quantum encoding slot.}
% Block A defines \label{sec:paper_v_l2_results}, which Blocks B and C cite.
%
% Macros used are already defined in the manuscript preamble: \tentative,
% \GKBA is unused here, \EOM is used as the manuscript defines it.  booktabs
% and revtex's ruledtabular are already available.
% ============================================================================


% ============================================================================
% BLOCK A --- new section, inserted immediately before \section{Numerical plan}
% ============================================================================

\section{Classical \(L=2\) diagnostics}
\label{sec:paper_v_l2_results}

The classical \(L=2\) solve called for in Sec.~\ref{sec:baseline_needed} has
been carried out for the driven Hubbard--Holstein dimer at strong coupling
\(\lambda=\tentative{1.5}\), \(\gamma=\tentative{0.5}\), drive amplitude
\(\tentative{1}\), phonon cutoff \(\tentative{16}\), from an exact correlated
preparation, with SSPRK(3,3) at \(h=\tentative{0.0025}\).  The reference is
the exact cutoff-\(\tentative{16}\) driven trajectory (DOP853,
\(\mathrm{rtol}=\tentative{10^{-9}}\), \(\mathrm{atol}=\tentative{10^{-11}}\),
maximum step \(\tentative{0.01}\)), which is a reporting input only and never
enters any propagation.  The error measure reported here is the
all-coordinate scalar RMS of the \(\tentative{31}\) closed scalar coordinates
against that reference, \(\sqrt{\operatorname{mean}(d^{2})}\) over samples and
coordinates on the \(h=\tentative{0.0025}\) grid, cumulative from \(t=0\).

The propagated object is a fixed-dictionary moment closure of
Eqs.~\eqref{eq:paper_v_rho_electron}--\eqref{eq:paper_v_delta_rho12q}, carried
together with the classical coherent field of Eq.~\eqref{eq:paper_v_B}.  At
its largest declared dictionary it has \(\tentative{27}\) entrance keys,
\(\tentative{56}\) propagated state coordinates and \(\tentative{36}\)
positive-extension rows over \(\tentative{70}\) frontier keys.  Physicality is
enforced by projection onto two cones --- the retained joint Gram of the raw
moment coordinates and the positive extension --- together with the hyperplane
on which the internal energy of Eq.~(22) of Ref.~\cite{RivaSimoniPing2026}
equals a reference advanced by the closure's own drive power.
Table~\ref{tab:paper_v_l2_scores} collects the result at the horizons the
record carries, beside the previously recorded routes and beside the
uncorrected archive equations of motion.

% ---------------------------------------------------------------------------
% Source for every entry of the table: the results summary, section 3.1 (the
% scores and the bands) and section 3.2 (the wall times and the machine note).
% Scores are the all-coordinate scalar RMS of the 31 closed scalar
% coordinates; the t = 0 coordinate offset of every route in the table is
% 8.08e-16, so raw and offset-subtracted figures agree to every printed digit.
% ---------------------------------------------------------------------------
\begin{table*}
\caption{Classical \(L=2\) trajectory error at strong coupling, in units of
\(10^{-2}\): all-coordinate scalar RMS of the \(31\) closed scalar
coordinates against the exact cutoff-\(16\) driven reference, cumulative from
\(t=0\) to each stated horizon.  The second line of a cell gives the number of
members of the measured band and its spread in percent; a band is the spread
of the run and its one-setting twins and is a measured sensitivity, not a
statistical error bar, so any ordering closer than the band is unresolved.
``1 run'' marks an entry with no band.  Common configuration:
\(\lambda=1.5\), \(\gamma=0.5\), drive amplitude \(1\), phonon cutoff \(16\),
exact correlated preparation, SSPRK(3,3) at \(h=0.0025\), lift re-selection
every \(16\) accepted steps.  The exact reference itself propagates to
\(t=4\) in \(0.7\,\mathrm{s}\).  Every figure is \tentative{tentative}: none
of it is adopted evidence.}
\label{tab:paper_v_l2_scores}
\begin{ruledtabular}
\footnotesize
\begin{tabular}{p{1.45in}cccccc}
configuration
 & \(t=4\) & \(t=8\) & \(t=12\) & \(t=16\) & \(t=20\)
 & wall (s) \\
\colrule
\(27\)-key dictionary, anchor-free measured re-selection rule, no trust
  radius, \(120\) Dykstra iterations, from \(t=0\)
 & \shortstack{\tentative{\(0.9290\)}\\{\scriptsize\tentative{\((4;\,7.1)\)}}}
 & \shortstack{\tentative{\(3.304\)}\\{\scriptsize\tentative{\((4;\,9.8)\)}}}
 & \shortstack{\tentative{\(4.060\)}\\{\scriptsize\tentative{(1 run)}}}
 & \shortstack{\tentative{\(4.610\)}\\{\scriptsize\tentative{(1 run)}}}
 & \shortstack{\tentative{\(5.604\)}\\{\scriptsize\tentative{(1 run)}}}
 & \tentative{\(5{,}183.6\)} \\
the same rule with the ball trust radius \(0.1\), \(30\) Dykstra iterations
 & \shortstack{\tentative{\(1.043\)}\\{\scriptsize\tentative{\((4;\,3.9)\)}}}
 & \shortstack{\tentative{\(3.181\)}\\{\scriptsize\tentative{\((4;\,2.3)\)}}}
 & \shortstack{\tentative{\(3.797\)}\\{\scriptsize\tentative{\((4;\,2.9)\)}}}
 & \shortstack{\tentative{\(4.821\)}\\{\scriptsize\tentative{\((4;\,5.9)\)}}}
 & \shortstack{\tentative{\(6.353\)}\\{\scriptsize\tentative{\((4;\,9.7)\)}}}
 & \tentative{\(4{,}844.7\)} \\
\(27\)-key dictionary, earlier re-selection form, \(30\) Dykstra iterations
 & \shortstack{\tentative{\(1.143\)}\\{\scriptsize\tentative{\((2;\,13.1)\)}}}
 & \shortstack{\tentative{\(2.593\)}\\{\scriptsize\tentative{\((2;\,4.5)\)}}}
 & \shortstack{\tentative{\(4.251\)}\\{\scriptsize\tentative{\((2;\,1.9)\)}}}
 & \shortstack{\tentative{\(6.377\)}\\{\scriptsize\tentative{\((2;\,5.4)\)}}}
 & \shortstack{\tentative{\(8.519\)}\\{\scriptsize\tentative{\((2;\,6.7)\)}}}
 & \tentative{\(4{,}794.4\)} \\
the first row with the Eq.~(22) hyperplane removed from the projection
 & \shortstack{\tentative{\(1.206\)}\\{\scriptsize\tentative{\((2;\,10.3)\)}}}
 & \shortstack{\tentative{\(4.853\)}\\{\scriptsize\tentative{\((2;\,1.9)\)}}}
 & \shortstack{\tentative{\(11.43\)}\\{\scriptsize\tentative{(1 run)}}}
 & \shortstack{\tentative{\(21.48\)}\\{\scriptsize\tentative{(1 run)}}}
 & \shortstack{\tentative{\(35.64\)}\\{\scriptsize\tentative{(1 run)}}}
 & \tentative{\(3{,}514.0\)} \\
\(15\)-key dictionary, ball rule, \(30\) Dykstra iterations
 & \shortstack{\tentative{\(2.72\)}\\{\scriptsize\tentative{\((4;\,0.1)\)}}}
 & \shortstack{\tentative{\(3.85\)}\\{\scriptsize\tentative{\((4;\,0.2)\)}}}
 & \shortstack{\tentative{\(9.89\)}\\{\scriptsize\tentative{\((4;\,5.5)\)}}}
 & \shortstack{\tentative{\(13.2\)}\\{\scriptsize\tentative{\((4;\,19.3)\)}}}
 & \shortstack{\tentative{\(13.3\)}\\{\scriptsize\tentative{\((4;\,56.2)\)}}}
 & \tentative{\(762\)--\(847\)} \\
\colrule
\(60\)-coordinate fixed-dictionary closure, full cone and retained joint-Gram
  controller
 & \tentative{\(1.2580\)} & --- & --- & --- & ---
 & \tentative{\(3{,}399.6\)}\footnotemark[1] \\
\(60\)-coordinate fixed-dictionary closure, both cones
 & --- & --- & --- & ---
 & \tentative{\(12.8607\)}
 & \tentative{\(5{,}832.34\)}\footnotemark[1] \\
plain archive \(\EOM\), no physicality enforcement
 & \tentative{\(11.0890\)} & \tentative{\(13.8174\)} & --- & ---
 & \tentative{\(19.3196\)}
 & \tentative{under \(1\)} \\
\end{tabular}
\footnotetext[1]{Measured on a different machine from every other wall in this
column; the ratio between the two machines is unmeasured, so these two entries
are not comparable with the rest of the column.}
\end{ruledtabular}
\end{table*}

Read across the table, no single configuration leads at every horizon: the
first row leads at \(t=\tentative{4}\) and \(t=\tentative{20}\), the third at
\(t=\tentative{8}\), and the second at \(t=\tentative{12}\) and
\(t=\tentative{16}\).

\paragraph*{Internal energy.}
The internal energy of Eq.~(22) of Ref.~\cite{RivaSimoniPing2026} is an
exactly linear functional of the raw moment coordinates, and enforcing it as a
third projection holds it on the declared reference to
\tentative{\(1.78\times10^{-15}\)} over \tentative{\(23{,}802\)} constrained
projections to \(t=\tentative{20}\) with no refusals; the sample RMS of that
energy over \(t=0\) to \(\tentative{20}\) is then
\tentative{\(5.336\times10^{-3}\)} for the first row of
Table~\ref{tab:paper_v_l2_scores}, against
\tentative{\(6.345\times10^{-3}\)} for the third,
\tentative{\(1.371\)} for the row with the hyperplane removed,
\tentative{\(9.396\times10^{-4}\)} for the both-cones line and
\tentative{\(1.824\times10^{-2}\)} for the uncorrected archive \(\EOM\), that
column measuring the level at which the reference itself sits
(\tentative{\(5.614\times10^{-3}\)} above the exact post-pulse energy
\tentative{\(-3.571397982\)}) rather than a drift.

\paragraph*{Physicality.}
Over \(t=0\) to \(\tentative{20}\) the first row of
Table~\ref{tab:paper_v_l2_scores} reaches a minimum unshifted retained
joint-Gram eigenvalue of \tentative{\(-4.460\times10^{-4}\)} (at
\(t=\tentative{11.04}\)) and a minimum scaled positive-extension eigenvalue of
\tentative{\(-3.936\times10^{-5}\)} (at \(t=\tentative{9.4975}\)), the
shallowest pair of any constrained chain on record, while no chain carrying
the Eq.~(22) hyperplane holds either cone past \(t=\tentative{8}\), the one
chain that does hold both (\tentative{\(-1.549\times10^{-12}\)} and
\tentative{\(-1.350\times10^{-9}\)}, both at the scale of the arithmetic) being
the one with that hyperplane removed and the worst in the table from
\(t=\tentative{8}\) onward.

\paragraph*{Wall and machine.}
Every wall in Table~\ref{tab:paper_v_l2_scores} except the two marked entries
was measured on one 24\,GB workstation, one heavy process at a time, each
rollout capped at \(\tentative{4500}\,\mathrm{MB}\) with measured peaks of
\tentative{\(624\)} to \tentative{\(699\)}\,MiB; the two marked entries come
from a different machine and are not comparable with the others.  None of the
figures in this section is adopted evidence.


% ============================================================================
% BLOCK B --- replaces the \placeholder line under
%             \paragraph*{Regularization slot.}  in \section{Numerical plan}
% ============================================================================

\paragraph*{Regularization slot.}
Of the five candidate high-\(U\) failure modes, the classical \(L=2\) case
exhibits density-matrix positivity loss: the uncorrected system leaves the
joint moment cone while every other listed defect stays at or near the
arithmetic floor.  The evidence excluding the other four is as follows.
\emph{Energy drift is not an intrinsic defect of the closure.}  The drive-free
velocity of the closed system changes the internal energy of Eq.~(22) of
Ref.~\cite{RivaSimoniPing2026} by \tentative{\(4.35\times10^{-17}\)} per step,
and the drift against the reference advanced by the closure's own drive power
is \tentative{\(4.09\times10^{-16}\)} per step; the uncorrected archive
\(\EOM\) carries a sample RMS of that energy of
\tentative{\(2.759\times10^{-2}\)} at \(t=\tentative{4}\) and
\tentative{\(1.824\times10^{-2}\)} at \(t=\tentative{20}\), which does not
accumulate with the horizon.  \emph{Correlation blow-up is not observed.}  The
uncorrected trajectory stays bounded over the whole horizon, its
all-coordinate scalar RMS against the exact cutoff-\(\tentative{16}\) driven
reference reading \tentative{\(1.109\times10^{-1}\)},
\tentative{\(1.382\times10^{-1}\)} and \tentative{\(1.932\times10^{-1}\)} at
\(t=\tentative{4}\), \(\tentative{8}\) and \(\tentative{20}\), while over the
same horizon its minimum unshifted retained joint-Gram eigenvalue falls to
\tentative{\(-5.173\times10^{-1}\)} by \(t=\tentative{4}\): the amplitudes
remain finite and it is representability that is lost.  \emph{The source
equations are not the defect, and are not themselves stiff.}  On the exact
state the archive equations of motion reproduce the exact flow --- the
residual-subtracted form scores \tentative{\(5.46\times10^{-7}\)} over
\(t=0\) to \(\tentative{0.05}\), and the phonon-covariance equation, Eq.~(14b)
of Ref.~\cite{RivaSimoniPing2026}, here Eq.~\eqref{eq:paper_v_delta_rho},
integrates at the exact state to the exact change of its own coordinate over
\(t=\tentative{4}\) to \(\tentative{8}\),
\tentative{\(2.323\times10^{-2}\)}, with residual
\tentative{\(5.010\times10^{-4}\)} --- and the uncorrected system reaches
\(t=\tentative{20}\) under a fixed step \(h=\tentative{0.0025}\) in
\tentative{under one second}.  \emph{Self-consistency instability is not
observed}: the coupled loop through \(\tilde h(t)\) and
Eq.~\eqref{eq:paper_v_B} runs the whole horizon at strong coupling without an
integrator failure.

\emph{Stiffness is real, but it is a property of the regularization and not of
Eqs.~\eqref{eq:paper_v_rho_electron}--\eqref{eq:paper_v_delta_rho12q}.}  This
is what the solver-step rejection counts and the adaptive integrator's step
history settle.  Under an adaptive integrator a containment-corrected velocity
built to hold the joint moment cone took \tentative{\(46{,}034\)} accepted
steps and \tentative{\(650{,}243\)} right-hand-side evaluations in
\tentative{\(9{,}000\)} seconds and reached \(t=\tentative{0.30}\); its first
\tentative{\(554\)} accepted steps carry the trajectory to
\(t=\tentative{0.1898}\) at a mean accepted step of
\tentative{\(3.4\times10^{-4}\)} and the remaining \tentative{\(45{,}480\)}
steps carry it \tentative{\(0.151\)} further at a mean accepted step of
\tentative{\(3.3\times10^{-6}\)}, a hundredfold collapse, while the projection
count per right-hand-side evaluation is flat at \tentative{\(40.5\)}, so the
loss is in the accepted step size and not in the price of one correction.  The
same uncorrected velocity is not stiff.  What limits a long horizon is
therefore not the target ODE but the step rejection induced by whatever
enforces its representability, and a regularization for this system is to be
judged on the stiffness it adds.
Section~\ref{sec:paper_v_l2_results} reports what a fixed-step projection
route achieves once that is accepted.


% ============================================================================
% BLOCK C --- replaces the \placeholder line under
%             \paragraph*{Quantum encoding slot.}  in \section{Numerical plan}
% ============================================================================

\paragraph*{Quantum encoding slot.}
The classical diagnostics of Sec.~\ref{sec:paper_v_l2_results} identify the
dominant bottleneck as invariant preservation rather than state dimension,
stiffness of the target, nonlinearity or measurement burden, and they localize
it further: what costs is one semidefinite projection, acting through one
coordinate of the phonon sector.  Of the wall of the first row of
Table~\ref{tab:paper_v_l2_scores}, \tentative{\(0.803\)} is spent in the
\tentative{\(500\)} positive-extension completion solves that keep the lift
representable, at \tentative{\(8.32\)} seconds each; the same share is
\tentative{\(0.838\)} for the earlier projective route
(\tentative{\(2{,}706.8\)} of \tentative{\(3{,}229.8\)} seconds, on
\tentative{\(1{,}068\)} of \tentative{\(5{,}559\)} stages) and
\tentative{\(0.981\)} for the strictly guarded route
(\tentative{\(142.350\)} of \tentative{\(145.172\)} seconds in
\tentative{\(36\)} conic solves).  Evaluating
Eqs.~\eqref{eq:paper_v_rho_electron}--\eqref{eq:paper_v_delta_rho12q} and
Eq.~\eqref{eq:paper_v_B} without that enforcement reaches \(t=\tentative{20}\)
in \tentative{under one second}.  The residual error the route does carry is
concentrated in the same place: of the phonon-occupation error over
\(t=\tentative{4}\) to \(\tentative{8}\), \tentative{\(0.998\)} is the
relative-mode phonon population \(v^{\dagger}\delta\rho\,v\) with
\(v=(1,-1)/\sqrt{2}\) alone, and it is the positive-extension projection that
moves it (\tentative{\(1.113\times10^{-3}\)} per stage repair, against
\tentative{\(7.5\times10^{-8}\)} for the retained joint-Gram projection).  At
fixed configuration the price of that enforcement is measurable in both
directions: deepening the projection from \(\tentative{30}\) to
\(\tentative{120}\) alternations moves the minimum retained joint-Gram
eigenvalue from \tentative{\(-2.38\times10^{-3}\)} to
\tentative{\(-4.42\times10^{-4}\)} and the \(t=\tentative{20}\) score from
\tentative{\(6.353\times10^{-2}\)} to \tentative{\(5.784\times10^{-2}\)} for
\tentative{\(11\)} percent more wall, while removing the energy hyperplane
returns \tentative{\(32.2\)} percent of the wall and costs \tentative{\(6.36\)}
times the score and \tentative{\(257\)} times the internal-energy RMS at
\(t=\tentative{20}\).

This reorders the three routes of Sec.~\ref{sec:pedagogical_solver_map}
without yet deciding among them.  The NISQ residual layer becomes \emph{less}
attractive: its object is the projected propagation of \(x\), which is the
part already measured to cost \tentative{under one second} of a
\tentative{\(5{,}183.6\)}-second wall, and
Eqs.~\eqref{eq:paper_v_sector_scale_eq}
and~\eqref{eq:paper_v_sector_theta_eq} carry no mechanism that keeps the
represented moment state inside either cone.  The fault-tolerant linearization
layer becomes less attractive for the same reason, since block-encoding
\(J_Q\) accelerates the same cheap object.  The fault-tolerant lifting layer
is the one whose object is the right one, because the lifted vector \(X_K\)
\emph{is} a moment vector and the constraint that costs is a semidefinite
condition on its Gram; and that condition is measured to be exactly affine in
the propagated coordinates (reconstruction errors
\tentative{\(1.1\times10^{-16}\)}, \tentative{\(1.1\times10^{-16}\)} and
\tentative{\(2.8\times10^{-17}\)} for the retained joint Gram, the positive
extension and the unified Gram, and at most \tentative{\(1.8\times10^{-15}\)}
at independent probes four times larger in scale), so the repair is a fixed
linear map, a spectral clip, and one least-squares map back whose design has
full rank \tentative{\(126\)} over the \tentative{\(126\)} movable
coordinates.  That triple is a matrix-function and linear-system object, which
is what the LCU/QSP and quantum-linear-system items of
Sec.~\ref{sec:ft_ideas} address; it is not a linear-ODE object, and the
lifting layer's own cost model in Sec.~\ref{sec:pedagogical_solver_map} does
not yet carry it.

Two measurements would settle the choice, and neither has been made.  First,
how the wall divides between velocity evaluation and the semidefinite step as
the dictionary grows: the share is measured at one dictionary
(\tentative{\(0.803\)} at \tentative{\(27\)} entrance keys and
\tentative{\(36\)} extension rows), and the derivative map of the extension
has been measured to have rank \tentative{\(1\)} over the four raw phonon
normal coordinates at all \tentative{\(11\)} dictionaries examined
(\tentative{\(24\)} to \tentative{\(45\)} extension rows), so the scaling of
the two shares with dictionary size and with \(L\) is the open quantity.  At
\(L=2\) it cannot be read off, because \(\dim x\) is only
\tentative{\(56\)} propagated real coordinates and no encoding claim can rest
on a dimension that small.  Second, the conditioning of the repair itself: the
amplification factor of the least-squares map back, together with the number
of alternations the cone actually requires --- which is measured to matter,
the projection running \tentative{three} times per accepted step and
\tentative{\(23{,}802\)} times to \(t=\tentative{20}\).  Until both exist,
item~4 of the decision rule in Sec.~\ref{sec:pedagogical_solver_map} stands:
none of the three layers currently prices the operation that dominates the
classical cost, so no route may be selected on the strength of a cost model
that omits it.
